% scopiazzato dal template di Matteo Longeri (grazie!)
%%%%%%%%%%%%%%%%%%%%%%%%%%%%%%%%%%%%%%%%%%%%%%%%%%%%%%
\documentclass[12pt,a4paper]{report}
% o article, book, ...



%%%%%%%%%%%%%%%%%%%%%%%%%%%%%%%%%%%%%%%%%%%%%%%%%%%%%%
% packages...
\usepackage[utf8]{inputenc}
\usepackage[english,italian]{babel}
\usepackage[hyphens]{url}

% Per generare il file PDF aderente alle specifiche PDF/A-1b. Verificarne poi la validità.
%\usepackage[a-1b]{pdfx}

\usepackage{hyperref}
\usepackage{graphicx}

\usepackage{underscore} % !!!
\usepackage{hyperref}
%\usepackage{bera}% optional: just to have a nice mono-spaced font
\usepackage{listings}
\usepackage{xcolor}

\colorlet{punct}{red!60!black}
\definecolor{background}{HTML}{EEEEEE}
\definecolor{delim}{RGB}{20,105,176}
\colorlet{numb}{magenta!60!black}

\lstdefinelanguage{json}{
    basicstyle=\normalfont\ttfamily,
    numbers=left,
    numberstyle=\scriptsize,
    stepnumber=1,
    numbersep=8pt,
    showstringspaces=false,
    breaklines=true,
    frame=lines,
    backgroundcolor=\color{background},
    literate=
     *{0}{{{\color{numb}0}}}{1}
      {1}{{{\color{numb}1}}}{1}
      {2}{{{\color{numb}2}}}{1}
      {3}{{{\color{numb}3}}}{1}
      {4}{{{\color{numb}4}}}{1}
      {5}{{{\color{numb}5}}}{1}
      {6}{{{\color{numb}6}}}{1}
      {7}{{{\color{numb}7}}}{1}
      {8}{{{\color{numb}8}}}{1}
      {9}{{{\color{numb}9}}}{1}
      {:}{{{\color{punct}{:}}}}{1}
      {,}{{{\color{punct}{,}}}}{1}
      {\{}{{{\color{delim}{\{}}}}{1}
      {\}}{{{\color{delim}{\}}}}}{1}
      {[}{{{\color{delim}{[}}}}{1}
      {]}{{{\color{delim}{]}}}}{1},
}


% Per inserire testo a caso in attesa di realizzare i capitoli
\usepackage{lipsum}



%%%%%%%%%%%%%%%%%%%%%%%%%%%%%%%%%%%%%%%%%%%%%%%%%%%%%
\begin{document}

% Frontespizio
\begin{titlepage}
  \begin{center}
    \includegraphics[width=\textwidth]{Logo.jpg}\\
    {\large{\bf Corso di Laurea in Informatica}}
  \end{center}
  \vspace{12mm}
  \begin{center}
    {\huge{\bf CONVERSIONE DI}}\\
    \vspace{4mm}
    {\huge{\bf STRUMENTI VINTAGE}}\\
    \vspace{4mm}
    {\huge{\bf PER LA DATA}}\\
    \vspace{4mm}
    {\huge{\bf PHYSICALIZATION}}\\
  \end{center}
  \vspace{12mm}
  \begin{flushleft}
    {\large{\bf Relatore:}}
    {\large{Andrea Trentini}}\\
    %\vspace{4mm}
    %{\large{\bf Correlatore:}}
    %{\large{...}}\\
  \end{flushleft}
  \vspace{12mm}
  \begin{flushright}
    {\large{\bf Tesi di Laurea di:}}
    {\large{Davide Busolin}}\\
    {\large{\bf Matr. 930814}}\\
  \end{flushright}
  \vspace{4mm}
  \begin{center}
    {\large{\bf Anno Accademico 2020/2021}}
  \end{center}
\end{titlepage}


\tableofcontents


% REPORT
%\part
%	\chapter
%		\section
%			\subsection
%				\paragraph
%					\subparagraph

% o sections (dipende dal documentclass)
\chapter{Introduzione}
TCP/IP over Avian Carriers\cite{waitzman1990standard}

\chapter{Attori coinvolti}
Prima di parlare del funzionamento dell'oggetto è il caso di soffermarsi su chi interagisce con esso:
\subsection*{Sviluppatore}
È in grado di apportare modifiche al codice sorgente. Può aggiungere, rimuovere e  modificare funzionalità del programma
e apportare modifiche dirette ai file di configurazione.
\subsection*{Utente esperto}
Conosce le API ed è in grado di comprendere il formato dei dati che restituiscono. Gli è sufficiente un'interfaccia anche spartana
per selezionare la sorgente dei dati e il campo specifico che vuole rappresentato.
\subsection*{Utente inesperto}
Non conosce le API né il formato JSON: ha bisogno di un prodotto già pronto che richieda la minima configurazione possibile e questa
deve essere particolarmente intuitiva, ad esempio una piccola interfaccia web per scegliere la rete WiFi e inserirne la password
al primo avvio. La sorgente dei dati deve essere preconfigurata e se ne viene resa disponibile più di una la scelta deve essere molto
semplice, possibilmente tramite interazione fisica con il dispositivo.


\chapter{Funzionamento del prototipo}

\section{All'avvio}
Per prima cosa contestualmente all'accensione, mediante WiFiManager, viene effettuato un tentativo di connessione all'ultima rete WiFi
utilizzata, di cui sono state salvate le credenziali. Se non è mai stata effettuata la connessione ad alcuna rete WiFi o se
quella salvata non è disponibile, viene attivato un Access Point che tramite Captive Portal (lo stesso che viene usato su reti pubbliche
o aziendali per richiedere l'autenticazione) chiede di selezionare una delle reti WiFi vicine e di inserirne la password.\\

\vspace{4mm}
{[SCREENSHOT]}\\
\vspace{4mm}

Successivamente viene tentata la connessione. Se non ha successo viene riproposto lo stesso procedimento, altrimenti viene disattivato
l'Access Point temporaneo e il programma prosegue.

\pagebreak
\section{Modalità operative}

\subsection{Modalità API}
Questa modalità di operazione si basa sull'esecuzione ad intervalli di tempo configurabili di richieste HTTP o HTTPS ad API REST
e sulla conversione del risultato in una tensione di uscita.
La configurazione è salvata nella memoria flash del microcontrollore
in formato JSON, esempio:

\begin{lstlisting}[language=json,firstnumber=1]
{
  "apiUrl": "http://api.coindesk.com/v1/bpi/currentprice/USD.json",
  "filterJSON": "{bpi:{USD:{rate_float:true}}}",
  "path": "bpi/USD/rate_float",
  "min_value": 58700,
  "max_value": 58900,
  "min_pwm": 0,
  "max_pwm": 1023,
  "request_interval_ms": 15000
}
\end{lstlisting}


\noindent Dove:
\begin{itemize}
  \item \texttt{apiUrl} è l'URL della risorsa (stringa)
  \item \texttt{post_payload} in caso di richiesta POST contiene la stringa da inviare nel body
  \item \texttt{filterJSON} è un documento JSON che contiene \texttt{true} come placeholder del campo che si vuole considerare dalla risposta (stringa/oggetto)
  \item \texttt{path} è una stringa che contiene il "percorso" del campo che si vuole rappresentare fisicamente (campo: float, path: stringa)
  \item \texttt{min_value} indica il valore minimo della scala per rappresentare il valore (float)
  \item \texttt{max_value} indica il valore massimo (float)
  \item \texttt{min_pwm} indica il valore minimo di uscita della PWM da mappare al valore restituito dalle API (int)
  \item \texttt{max_pwm} indica il valore massimo [su ESP il duty cycle massimo equivale a 1023] (int)
  \item \texttt{request_interval_ms} indica l'intervallo di tempo in millisecondi che intercorre tra le richieste alla risorsa (int)
\end{itemize}

La configurazione può essere alterata tramite una pagina web accessibile all'indirizzo IP nella rete locale del microcontrollore o il suo
hostname.

\vspace{4mm}
{[SCREENSHOT PAGINA WEB]}

\vspace{4mm}
I form sono precompilati con la configurazione in esecuzione (\emph{running-conf}) ed è presente una sezione che riporta la configurazione
salvata nella flash (\emph{saved-conf}).\\
È possibile inoltre inviare un JSON di configurazione tramite MQTT al topic \texttt{Phys/setFromJSON}. I campi possono essere modificati
singolarmente ai topic:
\begin{itemize}
  \item \texttt{Phys/setApiUrl}
        % ToDo: setPostPayload
  \item \texttt{Phys/setFilterJson}
  \item \texttt{Phys/setPath}
  \item \texttt{Phys/setMinValue}
  \item \texttt{Phys/setMaxValue}
  \item \texttt{Phys/setMinPwm}
  \item \texttt{Phys/setMaxPwm}
  \item \texttt{Phys/setRequestIntervalMs}
\end{itemize}

Le modifiche non sono automaticamente salvate nella memoria flash ma ciò deve essere richiesto esplicitamente. Nei form di configurazione
è presente una checkbox da spuntare in caso si voglia che questo avvenga.


\subsection{Modalità MQTT}
Non ancora implementata, ma l'idea è che consenta di visualizzare un valore ricevuto direttamente tramite sottoscrizione a un topic MQTT
invece che andarlo a reperire tramite richieste HTTP(S).

\chapter{Realizzazione del prototipo}

\section{Hardware}
Il prototipo è molto semplice dal punto di vista hardware, consiste in un voltmetro analogico ad ago e un WeMos D1 mini basato su ESP8266.\\
Il terminale positivo del voltmetro è connesso a un GPIO dell'ESP, mentre il terminale negativo a terra. È configurato mediante i
commutatori in modalità corrente continua con fondo scala a 3V.

\section{Software}
La parte software del prototipo è stata realizzata utilizzando il linguaggio di programmazione di Arduino, esteso all'utilizzo con
ESP8266 e con l'aiuto di numerose librerie.

\subsection{PWM}
Il principio di base che permette di produrre una tensione variabile da parte di un sistema digitale è la Pulse-Width Modulation.\\
Ciò permette di controllare finemente la posizione dell'ago del voltmetro via software.

\subsection{Richieste API}
Il primo passo verso la fisicalizzazione dei dati è stato quello di ottenerli, mediante richieste ad API REST in locale o su Internet.
Per questo è stata utilizzata la libreria \texttt{ESP8266HTTPClient}.\\
Le richieste HTTP hanno funzionato fin da subito, per HTTPS c'è stato più da lavorare in quanto non erano eccessivamente
documentate e tutte le soluzioni trovate online prevedevano l'utilizzo di stringhe per memorizzare intere risposte, che per via
della limitata disponibilità di memoria dell'ESP non è sempre una soluzione funzionante.\\
Un'alternativa fornita dalla libreria è quella di utilizzo di uno stream, per poter leggere dalla risposta un byte alla volta senza doverla memorizzare per intero, potendo ignorare le parti non necessarie. Questo meccanismo però è di facile utilizzo solo finché ci si limita a
HTTP. Per HTTPS la funzione per estrarre lo stream della risposta, secondo la documentazione della libreria, non è supportata.
Come soluzione è stato utilizzato come stream direttamente il WiFiClient utilizzato dalla connessione, che però non sempre presenta
 un output "pulito" della risposta ma potrebbe essere preceduto da una linea contentente (presumibilmente) la lunghezza in byte dell
 risposta.

\subsection{JSON}
La quasi totalità delle API mette a disposizione risposte contenenti dati in formato JSON (JavaScript Object Notation), è stata quindi
utilizzata la libreria \texttt{ArduinoJson} per la deserializzazione e la scelta del campo da estrarre.\\
JSON è stato successivamente usato anche come formato per la serializzazione e deserializzazione delle impostazioni da salvare sulla
memoria flash dell'ESP per permetterne la persistenza anche a seguito di perdite di tensione.

\subsection{Map}
Una volta ottenuto il valore da rappresentare è necessario trasformarlo in una tensione, ottenuta tramite PWM il cui duty cycle è
proporzionale a un numero a 10 bit (0 - 1023). Per fare questo è possibile utilizzare la funzione \texttt{map} già presente nella
libreria standard di Arduino.\\ 
Durante i test è emerso che tale funzione non si comporta particolarmente bene quando uno dei range è formato da piccoli numeri
decimali. Si è quindi preferito inserire nel codice direttamente la formula matematica su cui si basa \texttt{map}:
\[out = (value - fromLow) \times (toHigh - toLow) / (fromHigh - fromLow) + toLow \]

\noindent Dove:\cite{arduinomap}
\begin{itemize}
	\item value: Il numero da mappare.
	\item fromLow: Il limite inferiore dell'intervallo corrente del valore
	\item fromHigh: Il limite superiore dell'intervallo corrente del valore
	\item toLow: Il limite inferiore dell'intervallo di destinazione del valore
	\item toHigh: Il limite superiore dell'intervallo di destinazione del valore
\end{itemize}

\subsection{Task}
Attività periodiche come le richieste API sono state rese dei task per mezzo della libreria \texttt{TaskScheduler}, che permette di
impostare un intervallo di tempo tra esecuzioni ripetute di una procedura.

\subsection{Configurazione via MQTT}
Come prima modalità di modifica "a runtime" della configurazione, è stata implementata la possibilità di inviare messaggi MQTT ai
topic precedentemente descritti

\subsection{Credenziali WiFi}
Per la connessione alla rete wireless, grazie alla libreria \texttt{WiFiManager}, l'ESP memorizzerà le ultime credenziali utilizzate
per ritentare la connessione al riavvio. In caso le credenziali non siano valide o non ve ne siano di salvate, verrà avviato un 
Access Point temporaneo con captive portal per richiedere delle credenziali valide.

\subsection{Configurazione tramite form web}
Un'altra forma di modifica della configurazione è stata implementata tramite un web server configurato per mezzo della libreria
\texttt{ESPAsyncWebServer}. Inizialmente la pagina conteneva un semplice form tramite il quale era possibile inviare un intero JSON 
di configurazione, successivamente è stata aggiunta la possibilità di impostare i singoli parametri attraverso un form precompilato
e richieste AJAX. Un punto debole di questa soluzione è la scarsa flessibilità del web server, che gestisce richieste POST solo codificate
come Form Data, quindi invece che poter mandare un semplice JSON nel body della richiesta come da prassi con AJAX, è stato necessario
trasformarlo in FormData come valore associato alla chiave setFromJSON.\\
Nel dettaglio, sono esposti i seguenti endpoint:
\begin{itemize}
  \item /: Pagina di configurazione (index.html)
  \item /get/running-conf: Restituisce la configurazione attualmente in esecuzione in formato JSON
  \item /get/saved-conf: Restituisce il file di configurazione presente nella memoria flash
  \item /post: permette di inviare un intero JSON di configurazione (anche incompleto, vengono modificati solo i campi specificati) come
              valore associato alla chiave setFromJSON. È inoltre possibile asssociare alla chiave saveToFlash il valore `true' per copiare
              la configurazione in esecuzione nella memoria flash.
\end{itemize}
%%%%%% Da rivedere

\chapter{Conclusioni}

\section{Reperibilità delle API}
Durante la ricerca di API rappresentabili mediante Data Physicalization è risultato evidente che molte di quelle più congeniali non
sono pubbliche e spesso sono anche di difficile utilizzo. L'esperienza personale è stata con quelle di ATM per ottenere i minuti rimanenti
all'arrivo di un mezzo pubblico ma altri si sono scontrati con quelle di Trenitalia \cite{trenitaliashock}

\subsection{Dettagli su ATM}
Le API di ATM sono usate per \href{https://giromilano.atm.it}{Giromilano} e l'unico modo per comprenderne il funzionamento è utilizzare
gli strumenti del browser per leggere le richieste e le risposte che vengono scambiate tra client e server durante la ricerca di un
percorso o di informazioni su una fermata specifica.\\
In particolare per ottenere i minuti rimanenti all'arrivo di un mezzo ad una specifica fermata, è necessario effettuare una richiesta
POST a\\ \url{https://giromilano.atm.it/proxy.ashx} e nel body inserire\\ \texttt{url=tpPortal/geodata/pois/stops/xxxxx}, dove a
\texttt{xxxxx} va sostituito il numero della fermata, leggibile (a volte) nel mondo reale sulla pensilina o sul cartello o direttamente
dal sito.\\
Per ottenere una risposta è inoltre necessario inviare l'header\\ \texttt{Content-Type: application/x-www-form-urlencoded}.\\
A questo punto sorgono alcuni problemi:
\begin{itemize}
  \item Le risposte contengono un po' di tutto e per via degli avvisi comprensivi di tag HTML possono anche essere molto lunghe, anche
        decine di kb.
  \item In caso di richiesta non corretta, ad esempio mancante di header Content-Type, non si avranno in risposta codici HTTP specifici
        ma sempre 200 (OK) in cui però la risposta è vuota.
\end{itemize}

A prescindere dalle questioni tecniche il problema è anche legale: non essendo API pubbliche l'utilizzo che se ne può fare all'esterno
dell'ambito aziendale è molto limitato. Può andare bene per utilizzi sperimentali personali, ma non si possono certamente incorporare
in prodotti da distribuire a qualsiasi titolo.\\
Poter risalire al ritardo accumulato da un mezzo pubblico (specialmente per treni con cadenza oraria o semioraria) senza dover navigare
siti e app ma direttamente lanciando un'occhiata a un angolo della stanza, come si farebbe con un orologio, sarebbe una comodità in più
per i pendolari che al momento non è attuabile su larga scala da terzi e le aziende di trasporti non sembrano interessate a proporre
prodotti simili.



\bibliographystyle{plain}
\bibliography{Biblio}
%\addcontentsline{toc}{chapter}{Bibliografia}

\end{document}
